\documentclass[12pt]{article}
\usepackage[ukrainian]{babel}
\usepackage{fontspec}
% --- НАЛАШТУВАННЯ ШРИФТІВ ---
\setmainfont{Schoolbook-Regular.otf}[
    Path = ../,
    BoldFont = Schoolbook-Bold.otf,
    ItalicFont = Schoolbook-Italic.otf,
    BoldItalicFont = Schoolbook-BoldItalic.otf
]
\usepackage[italic]{mathastext}
\usepackage[a4paper,margin=1.8cm,bottom=2cm,top=2cm]{geometry}
\usepackage{amsmath,amssymb}
\usepackage{tikz}
\usepackage{pgfplots}
\pgfplotsset{compat=1.16}
\usetikzlibrary{calc,patterns,angles,quotes,intersections}
\usepackage{xcolor,array,fancyhdr,enumitem,multicol}
\usepackage{colortbl}
\usepackage{tcolorbox}
\tcbuselibrary{skins,breakable}
\usepackage[hidelinks]{hyperref}
\usepackage[firstpage=false, text=@pvtr2525, scale=0.6, color=gray!12]{draftwatermark}

\definecolor{mainGreen}{RGB}{34, 120, 64}
\definecolor{yearOrange}{RGB}{220, 100, 30}
\definecolor{yearcolor}{RGB}{220, 100, 30}
\definecolor{titleBg}{RGB}{225, 240, 220}
\definecolor{instrBg}{RGB}{255, 240, 220}
\definecolor{instrBorder}{RGB}{220, 150, 80}

\pagestyle{fancy}
\fancyhf{}
\renewcommand{\headrulewidth}{0.4pt}
\renewcommand{\footrulewidth}{0.4pt}
\fancyhead[C]{\small\color{gray!80} Збірник НМТ \textendash{} Варіант 8 червня 2026}
\fancyhead[R]{\small\color{mainGreen}\textbf{\href{https://t.me/pvtr2525}{@pvtr2525}}}
\fancyfoot[L]{\small\color{gray!80}\href{https://t.me/pvtr2525}{ТГ: @pvtr2525}}
\fancyfoot[C]{\small\color{gray!80}\thepage}
\fancyfoot[R]{\small\color{gray!80}НМТ 2026}

% --- ТАБЛИЦІ ВІДПОВІДЕЙ ---
\newcommand{\answerTable}[5]{%
\par\vspace{0.15cm}
\noindent
\begin{tabular}{|*{5}{>{\centering\arraybackslash}m{3cm}|}}
\hline
\rule[-0.15cm]{0pt}{0.6cm}\textbf{А} & \rule[-0.15cm]{0pt}{0.6cm}\textbf{Б} & \rule[-0.15cm]{0pt}{0.6cm}\textbf{В} & \rule[-0.15cm]{0pt}{0.6cm}\textbf{Г} & \rule[-0.15cm]{0pt}{0.6cm}\textbf{Д} \\
\hline
\rule[-0.2cm]{0pt}{0.75cm}#1 & \rule[-0.2cm]{0pt}{0.75cm}#2 & \rule[-0.2cm]{0pt}{0.75cm}#3 & \rule[-0.2cm]{0pt}{0.75cm}#4 & \rule[-0.2cm]{0pt}{0.75cm}#5 \\
\hline
\end{tabular}
\par\vspace{0.25cm}
}
\newcommand{\answerTableTall}[5]{%
\par\vspace{0.15cm}
\noindent
\begin{tabular}{|*{5}{>{\centering\arraybackslash}m{3cm}|}}
\hline
\rule[-0.2cm]{0pt}{0.7cm}\textbf{А} & \rule[-0.2cm]{0pt}{0.7cm}\textbf{Б} & \rule[-0.2cm]{0pt}{0.7cm}\textbf{В} & \rule[-0.2cm]{0pt}{0.7cm}\textbf{Г} & \rule[-0.2cm]{0pt}{0.7cm}\textbf{Д} \\
\hline
\rule[-0.45cm]{0pt}{1.2cm}#1 & \rule[-0.45cm]{0pt}{1.2cm}#2 & \rule[-0.45cm]{0pt}{1.2cm}#3 & \rule[-0.45cm]{0pt}{1.2cm}#4 & \rule[-0.45cm]{0pt}{1.2cm}#5 \\
\hline
\end{tabular}
\par\vspace{0.25cm}
}

\newcommand{\instructionBox}[1]{%
\par\vspace{0.3cm}
\begin{tcolorbox}[colback=instrBg, colframe=instrBorder, boxrule=0.6pt, arc=2pt,
    left=8pt, right=8pt, top=4pt, bottom=4pt]
\centering\small\bfseries #1
\end{tcolorbox}
\par\vspace{0.15cm}
}

\newcommand{\sectionTitle}[1]{%
\par\vspace{0.4cm}
\begin{tcolorbox}[colback=titleBg, colframe=titleBg, boxrule=0pt, arc=8pt,
    halign=center, fontupper=\Large\bfseries\color{mainGreen}]
#1
\end{tcolorbox}
\par\vspace{0.3cm}
}

\newcommand{\nmtAnswerBox}{%
\par\vspace{0.2cm}\noindent
Відповідь:\ \
\begin{tabular}{@{}*{4}{|>{\centering\arraybackslash}p{0.8cm}}|@{\hspace{0.1cm}\textbf{,}\hspace{0.1cm}}*{3}{|>{\centering\arraybackslash}p{0.8cm}}|@{}}
\hline
\rule[-0.2cm]{0pt}{0.7cm} & & & & & & \\
\hline
\end{tabular}
\par\vspace{0.3cm}
}

\newcommand{\matchingGrid}{%
    \begingroup
    \renewcommand{\arraystretch}{1.15}
    \setlength{\tabcolsep}{4pt}
    \footnotesize
    \begin{tabular}{r|c|c|c|c|c|}
         \multicolumn{1}{c}{} & \multicolumn{1}{c}{\textbf{А}} & \multicolumn{1}{c}{\textbf{Б}} & \multicolumn{1}{c}{\textbf{В}} & \multicolumn{1}{c}{\textbf{Г}} & \multicolumn{1}{c}{\textbf{Д}} \\ \cline{2-6}
         \textbf{1} & \phantom{X} & \phantom{X} & \phantom{X} & \phantom{X} & \phantom{X} \\ \cline{2-6}
         \textbf{2} & & & & & \\ \cline{2-6}
         \textbf{3} & & & & & \\ \cline{2-6}
    \end{tabular}
    \endgroup
}

\newcommand{\taskBlock}[1]{%
\par\vspace{0.45cm}
\noindent{\color{mainGreen}\rule{\linewidth}{0.8pt}}\par\vspace{0.12cm}
\noindent{\small\bfseries\color{mainGreen}#1}\par\vspace{0.25cm}
}

\newcounter{zad}
\newcommand{\zadtask}[1]{\stepcounter{zad}\noindent\textbf{\thezad.}\ \ #1\par\vspace{0.2cm}}
\newcommand{\zadnum}{\stepcounter{zad}\textbf{\thezad.}\ \ }

\setcounter{tocdepth}{2}
\newcommand{\chapterTitle}[1]{%
\clearpage\phantomsection\addcontentsline{toc}{section}{#1}%
\sectionTitle{#1}}

\begin{document}
\thispagestyle{empty}
\vspace*{1.2cm}
\begin{center}
{\Large\bfseries\color{mainGreen} ОРИГІНАЛЬНІ ЗАВДАННЯ НМТ-2026}\\[0.4cm]
{\Huge\bfseries\color{mainGreen} ВАРІАНТ ЗА 8 ЧЕРВНЯ}\\[0.6cm]
{\large Детальний розбір реального тесту з усіма рисунками та відповідями}\\[1cm]
{\large\color{mainGreen}\textbf{\href{https://t.me/pvtr2525}{Телеграм: @pvtr2525}}}
\end{center}
\clearpage

\chapterTitle{ТЕСТОВИЙ ЗОШИТ (8 ЧЕРВНЯ)}
\instructionBox{Завдання 1--15 мають по п'ять варіантів відповіді, з яких лише один правильний. Виберіть правильний варіант відповіді.}

\zadtask{Пляшку об'ємом $750$~мл на $\dfrac{2}{3}$ заповнили соком. Скільки соку залишилося в ній, якщо відлили $150$~мл соку?}
\answerTable{$250$~мл}{$600$~мл}{$300$~мл}{$450$~мл}{$350$~мл}

\noindent
\begin{minipage}[c]{0.48\textwidth}
\zadnum На рисунку зображено графік залежності кількості переглядів відеоролика (у тисячах) від часу (у годинах) з моменту його публікації. Укажіть проміжок часу, протягом якого кількість переглядів зростала найшвидше.
\end{minipage}%
\hfill
\begin{minipage}[c]{0.50\textwidth}
\centering
\begin{tikzpicture}
\begin{axis}[
    width=8.5cm, height=5.5cm,
    xmin=0, xmax=21,
    ymin=0, ymax=24,
    xtick={0,2,4,6,8,10,12,14,16,18,20},
    ytick={0,4,8,12,16,20,24},
    xlabel={\small Години},
    ylabel={\small Кількість переглядів, тис.},
    grid=both,
    grid style={dashed, gray!40},
    axis lines=left,
    tick label style={font=\scriptsize},
    label style={font=\scriptsize}
]
\addplot[mark=*, thick, red!80!black] coordinates {
    (0,0) (2,1) (4,2) (6,3) (8,6) (10,12) (12,15) (14,16) (16,17) (18,17.5) (20,18)
};
\end{axis}
\end{tikzpicture}
\end{minipage}
\par\vspace{0.2cm}
\answerTable{$[0; 4]$}{$[4; 8]$}{$[8; 10]$}{$[12; 16]$}{$[18; 20]$}

\noindent
\begin{minipage}[c]{0.60\textwidth}
\zadnum Ділянка складається з прямокутного трикутника та прямокутника (див. рисунок). За наведеними на рисунку розмірами сторін визначте периметр цієї ділянки.
\end{minipage}%
\hfill
\begin{minipage}[c]{0.38\textwidth}
\centering
\begin{tikzpicture}[scale=0.3, thick, line join=round]
    \coordinate (A) at (0,0);
    \coordinate (D) at (10,0);
    \coordinate (C) at (10,7);
    \coordinate (B) at (0,7);
    \coordinate (Top) at (6.4, 11.8);
    \draw (A) -- (D) -- (C) -- (Top) -- (B) -- cycle;
    \draw[dashed] (B) -- (C);
    \pic[draw, angle radius=0.3cm] {right angle = A--B--C};
    \pic[draw, angle radius=0.3cm] {right angle = B--A--D};
    \pic[draw, angle radius=0.3cm] {right angle = A--D--C};
    \pic[draw, angle radius=0.3cm] {right angle = D--C--B};
    \pic[draw, angle radius=0.3cm] {right angle = B--Top--C};
    \node[left] at (0,3.5) {$7$ м};
    \node[right] at (10,3.5) {$7$ м};
    \node[above left] at (3.2, 9.4) {$8$ м};
    \node[above right] at (8.2, 9.4) {$6$ м};
\end{tikzpicture}
\end{minipage}
\par\vspace{0.2cm}
\answerTable{$28$~м}{$35$~м}{$48$~м}{$38$~м}{$94$~м}

\zadtask{Позначте число, що є розв'язком системи нерівностей $\begin{cases}x+10>0,\\ 5-x>0.\end{cases}$}
\answerTable{$-5$}{$-10$}{$10$}{$15$}{$5$}

\noindent
\begin{minipage}[c]{0.62\textwidth}
\zadnum На рисунку зображено правильну чотирикутну піраміду $SABCD$, $SO$~--- її висота. Скільки всього є прямих, перпендикулярних до $SO$ та які проходять через дві вершини основи піраміди?
\vspace{0.2cm}
\par\noindent\textbf{А}\ \ лише одна
\par\noindent\textbf{Б}\ \ дві
\par\noindent\textbf{В}\ \ чотири
\par\noindent\textbf{Г}\ \ шість
\par\noindent\textbf{Д}\ \ більше шести
\end{minipage}%
\hfill
\begin{minipage}[c]{0.36\textwidth}
\centering
\begin{tikzpicture}[scale=0.8, thick, line join=round]
    \coordinate (A) at (0,0);
    \coordinate (B) at (1.5,1.2);
    \coordinate (C) at (5,1.2);
    \coordinate (D) at (3.5,0);
    \coordinate (O) at (2.5,0.6);
    \coordinate (S) at (2.5,4.5);
    \draw (A) -- (D) -- (C) -- (S) -- cycle;
    \draw (S) -- (D);
    \draw[dashed] (A) -- (B) -- (C);
    \draw[dashed] (S) -- (B);
    \draw[dashed] (A) -- (C);
    \draw[dashed] (B) -- (D);
    \draw[dashed] (S) -- (O);
    \node[below left] at (A) {$A$};
    \node[above left] at (B) {$B$};
    \node[right] at (C) {$C$};
    \node[below right] at (D) {$D$};
    \node[below] at (O) {$O$};
    \node[above] at (S) {$S$};
\end{tikzpicture}
\end{minipage}
\par\vspace{0.3cm}

\zadtask{$\dfrac{a^2-4b^2}{3a+6b} =$}
\answerTableTall{$\dfrac{a+2b}{2}$}{$\dfrac{a-2b}{3}$}{$3a-6b$}{$\dfrac{3}{a-2b}$}{$3(a+2b)$}

\zadtask{Якому проміжку належить корінь рівняння $2^{2x-5} = 8$?}
\answerTable{$(-\infty; -2)$}{$[0; 2)$}{$[-2; 0)$}{$[2; 4)$}{$[4; +\infty)$}

\noindent
\begin{minipage}[c]{0.58\textwidth}
\zadnum У прямокутній системі координат на площині зображено точки $K$, $L$, $M$, $N$ і $P$. Через яку із цих точок проходить графік функції $y=3$?
\end{minipage}%
\hfill
\begin{minipage}[c]{0.40\textwidth}
\centering
\begin{tikzpicture}[scale=0.6, thick]
    \draw[step=1cm, gray!40, thin] (-4,-3) grid (5,5);
    \draw[->, thick] (-4.5,0) -- (5.5,0) node[right] {$x$};
    \draw[->, thick] (0,-3.5) -- (0,5.5) node[above] {$y$};
    \node[below left] at (0,0) {$0$};
    \node[below right] at (1,0) {$1$};
    \node[above left] at (0,1) {$1$};
    
    \fill (-3,2) circle (3pt) node[above left] {$M$};
    \fill (1,3) circle (3pt) node[above left] {$N$};
    \fill (4,2) circle (3pt) node[above left] {$L$};
    \fill (0,-2) circle (3pt) node[below left] {$K$};
    \fill (3,-3) circle (3pt) node[above right] {$P$};
\end{tikzpicture}
\end{minipage}
\par\vspace{0.2cm}
\answerTable{$N$}{$K$}{$L$}{$M$}{$P$}

\zadtask{Які з наведених тверджень є правильними?
\vspace{0.1cm}
\begin{enumerate}[label=\textbf{\Roman*.}, align=left, leftmargin=2.6em, labelsep=0.5em, itemsep=0.08cm, topsep=0.06cm]
\item У будь-який ромб можна вписати коло.
\item У будь-яку прямокутну трапецію можна вписати коло.
\item У будь-якому трикутнику бісектриси кутів проходять через центр вписаного в цей трикутник кола.
\end{enumerate}
\vspace{0.15cm}}
\answerTable{лише III}{лише I та II}{лише I та III}{лише II та III}{I, II та III}
\newpage
\zadtask{Визначте всі критичні точки функції $f(x) = \dfrac{x^3}{3} - 2x^2 + 3x + 1$.}
\answerTable{$-3; -1$}{$1; 3$}{$3$}{$-3; 1$}{$1$}

\zadtask{Комп'ютерна програма видаляє у п'ятицифровому числі одну цифру навмання. Яка ймовірність того, що в числі $37281$ буде видалено цифру $1$ або цифру $2$?}
\answerTableTall{$\dfrac{1}{2}$}{$\dfrac{2}{5}$}{$\dfrac{1}{5}$}{$\dfrac{2}{3}$}{$\dfrac{1}{3}$}

\noindent
\begin{minipage}[c]{0.60\textwidth}
\zadnum У прямокутній системі координат у просторі зображено сферу радіуса $4$ (див. рисунок). Точки $M(1; -2; 0)$ і $N$ належать діаметру сфери. Визначте координати точки $N$.
\end{minipage}%
\hfill
\begin{minipage}[c]{0.38\textwidth}
\centering
\begin{tikzpicture}[scale=0.7, thick, line join=round]
    % Площина
    \draw[fill=yellow!30!orange!20, draw=black] (-2,-1) -- (4,-1) -- (5.5,1.5) -- (-0.5,1.5) -- cycle;
    % Осі
    \draw[->, >=stealth] (2.5,0.5) -- (1,-1.5) node[left] {$x$};
    \draw[->, >=stealth] (2.5,0.5) -- (5,0.5) node[below] {$y$};
    \draw[->, >=stealth] (2.5,0.5) -- (2.5,4.5) node[right] {$z$};
    % Сфера
    \coordinate (C) at (1, 2.5); % центр сфери
    \draw[fill=cyan!15, draw=blue!60!black] (C) circle (1.5);
    \draw[dashed, blue!60!black] (1, 2.5) ellipse (1.5 and 0.4); % екватор
    % Точки M і N та діаметр
    \coordinate (M) at (1, 1);
    \coordinate (N) at (1, 4);
    \draw[dashed, red!80!black] (M) -- (N);
    \fill (M) circle (2pt) node[left] {$M$};
    \fill (N) circle (2pt) node[left] {$N$};
\end{tikzpicture}
\end{minipage}
\par\vspace{0.2cm}
\answerTable{$(1; -2; 4)$}{$(1; -2; 8)$}{$(4; -2; 4)$}{$(5; -2; 0)$}{$(-1; 2; 8)$}

\zadtask{$\dfrac{(a^3)^4}{a^4 \cdot a^8} =$}
\answerTableTall{$1$}{$a^{\frac{1}{2}}$}{$a^6$}{$0$}{$-1$}

\noindent
\begin{minipage}[c]{0.58\textwidth}
\zadnum На стороні $BC$ прямокутника $ABCD$ вибрано точку $K$ так, що прямі $AK$ і $CD$ перетинаються в точці $M$ (див. рисунок). $MD = 12$~см, $MC : CD = 1 : 3$. Визначте площу прямокутника.
\end{minipage}%
\hfill
\begin{minipage}[c]{0.40\textwidth}
\centering
\begin{tikzpicture}[scale=0.3, thick, line join=round]
    \coordinate (A) at (0,0);
    \coordinate (D) at (15,0);
    \coordinate (C) at (15,5);
    \coordinate (B) at (0,5);
    \coordinate (K) at (10,5);
    \coordinate (M) at (15,7.5);
    \draw (A) -- (B) -- (C) -- (D) -- cycle;
    \draw (A) -- (M);
    \draw (D) -- (M);
    \fill (K) circle (3pt);
    \fill (M) circle (3pt);
    \node[left] at (A) {$A$};
    \node[above left] at (B) {$B$};
    \node[right] at (C) {$C$};
    \node[right] at (D) {$D$};
    \node[above left] at (K) {$K$};
    \node[above right] at (M) {$M$};
    \pic[draw, angle radius=0.8cm, pic text={$30^\circ$}, angle eccentricity=1.5, font=\scriptsize] {angle = D--A--M};
\end{tikzpicture}
\end{minipage}
\par\vspace{0.2cm}
\answerTable{$54\sqrt{3}$}{$81\sqrt{3}$}{$108\sqrt{3}$}{$216$}{$162$}

\zadtask{Розв'яжіть рівняння $\log_5 x + \log_5(4x) = \log_5 25$.}
\answerTableTall{$-\dfrac{5}{2}; \dfrac{5}{2}$}{$\dfrac{5}{2}$}{$\dfrac{2}{5}$}{$-5$}{$5$}

\instructionBox{У завданнях 16--18 до кожного з трьох рядків інформації, позначених цифрами, доберіть один правильний варіант, позначений буквою.}

\zadtask{Узгодьте вираз (1--3) із проміжком (А--Д), якому належить значення цього виразу.
\vspace{0.3cm}
\noindent
\begin{minipage}[t]{0.45\textwidth}
\textit{Вираз}\par\vspace{0.15cm}
\textbf{1} \quad $\log_{0{,}5} 2$\\[0.25cm]
\textbf{2} \quad $\sqrt{3^2 + (-2)^2}$\\[0.25cm]
\textbf{3} \quad $\left| 2 \cdot \dfrac{2}{7} - 2\dfrac{2}{7} \right|$
\end{minipage}
\hfill
\begin{minipage}[t]{0.3\textwidth}
\textit{Проміжок}\par\vspace{0.15cm}
\textbf{А} \quad $[-2; -1)$\\[0.15cm]
\textbf{Б} \quad $[-1; 1)$\\[0.15cm]
\textbf{В} \quad $[1; 2)$\\[0.15cm]
\textbf{Г} \quad $[2; 3)$\\[0.15cm]
\textbf{Д} \quad $[3; 4)$
\end{minipage}
\hfill
\begin{minipage}[t]{0.2\textwidth}
\vspace{0.4cm}
\begin{flushright}\matchingGrid\end{flushright}
\end{minipage}}

\zadtask{Доберіть до функції (1--3) її властивість (А--Д).
\vspace{0.3cm}
\noindent
\begin{minipage}[t]{0.45\textwidth}
\textit{Функція}\par\vspace{0.15cm}
\textbf{1} \quad $y = \lg x$\\[0.25cm]
\textbf{2} \quad $y = 0{,}2^x - 1$\\[0.25cm]
\textbf{3} \quad $y = x^2 - 4x + 5$
\end{minipage}
\hfill
\begin{minipage}[t]{0.3\textwidth}
\textit{Властивість}\par\vspace{0.15cm}
\textbf{А} \quad є зростаючою\\[0.15cm]
\textbf{Б} \quad є спадною\\[0.15cm]
\textbf{В} \quad не має нулів\\[0.15cm]
\textbf{Г} \quad має лише два нулі\\[0.15cm]
\textbf{Д} \quad є парною
\end{minipage}
\hfill
\begin{minipage}[t]{0.2\textwidth}
\vspace{0.4cm}
\begin{flushright}\matchingGrid\end{flushright}
\end{minipage}}

\zadtask{На рисунку зображено рівні паралелограми $ABCD$ і $AKMD$. $BD \perp AB$. $\angle BAD = 65^\circ$. Установіть відповідність між кутом (1--3) та його градусною мірою (А--Д).
\vspace{0.3cm}
\begin{center}
\begin{tikzpicture}[scale=0.7, thick, line join=round]
    \coordinate (A) at (0,0);
    \coordinate (D) at (5,0);
    \coordinate (B) at (1.5, 3.21); % 65 deg
    \coordinate (C) at (6.5, 3.21);
    \coordinate (K) at (1.5, -3.21);
    \coordinate (M) at (6.5, -3.21);
    
    \draw (-1,0) -- (7,0); % Line AD extends
    \draw (A) -- (B) -- (C) -- (D);
    \draw (A) -- (K) -- (M) -- (D);
    \draw (B) -- (D);
    
    \pic[draw, angle radius=0.5cm, pic text=$65^\circ$, angle eccentricity=1.7, font=\scriptsize] {angle = D--A--B};
    \pic[draw, angle radius=0.3cm] {right angle = A--B--D};
    
    \node[above left] at (A) {$A$};
    \node[above left] at (B) {$B$};
    \node[above right] at (C) {$C$};
    \node[above right] at (D) {$D$};
    \node[below left] at (K) {$K$};
    \node[below right] at (M) {$M$};
\end{tikzpicture}
\end{center}
\vspace{0.2cm}
\noindent
\begin{minipage}[t]{0.45\textwidth}
\textit{Кут}\par\vspace{0.15cm}
\textbf{1} \quad $\angle ABC$\\[0.25cm]
\textbf{2} \quad $\angle CDM$\\[0.25cm]
\textbf{3} \quad $\angle BDM$
\end{minipage}
\hfill
\begin{minipage}[t]{0.3\textwidth}
\textit{Градусна міра}\par\vspace{0.15cm}
\textbf{А} \quad $115^\circ$\\[0.15cm]
\textbf{Б} \quad $125^\circ$\\[0.15cm]
\textbf{В} \quad $130^\circ$\\[0.15cm]
\textbf{Г} \quad $140^\circ$\\[0.15cm]
\textbf{Д} \quad $150^\circ$
\end{minipage}
\hfill
\begin{minipage}[t]{0.2\textwidth}
\vspace{0.4cm}
\begin{flushright}\matchingGrid\end{flushright}
\end{minipage}}

\instructionBox{Розв'яжіть завдання 19--22. Одержані числові відповіді запишіть у спеціально відведеному місці.}

\zadtask{Знайдіть площу фігури, обмеженої графіком функції $y=3-\sqrt{x}$ та осями координат.}
\nmtAnswerBox

\zadtask{Михайло планував купити мобільний телефон, чохол до нього та карту пам'яті. Вартість телефона становить $8000$~грн, чохла~--- $600$~грн, а карти пам'яті~--- $800$~грн. У магазині проходить акція: за умови одночасної купівлі телефона та чохла покупець отримує знижку розміром $15\,\%$ від вартості телефона, знижку на чохол у розмірі п'ятої частини від його повної вартості, a карту пам'яті у подарунок. Яку суму грошей (у грн) заплатить Михайло за вибрані ним телефон, чохол та карту пам'яті, якщо скористається цією акцією?}
\nmtAnswerBox

\noindent
\begin{minipage}[c]{0.60\textwidth}
\zadnum На рисунку зображено правильну трикутну призму та циліндр. Об'єм циліндра дорівнює $80\pi$~см$^3$, а площа основи призми дорівнює $12\sqrt{3}$~см$^2$. Радіус основи циліндра дорівнює радіусу кола, вписаного в основу призми. Визначте об'єм (у см$^3$) призми, якщо висота призми відноситься до висоти циліндра як $\sqrt{3} : 1$.
\end{minipage}%
\hfill
\begin{minipage}[c]{0.38\textwidth}
\centering
\begin{tikzpicture}[scale=0.6, thick, line join=round]
    % Призма
    \coordinate (A1) at (0,0);
    \coordinate (B1) at (2,-1);
    \coordinate (C1) at (3,0.5);
    \coordinate (A2) at (0,4);
    \coordinate (B2) at (2,3);
    \coordinate (C2) at (3,4.5);
    
    \draw (A1) -- (B1) -- (B2) -- (A2) -- cycle;
    \draw (B1) -- (C1) -- (C2) -- (B2);
    \draw (A2) -- (C2);
    \draw[dashed] (A1) -- (C1);
    
    % Циліндр
    \begin{scope}[shift={(6,0)}]
        \draw[dashed] (-1,0) arc (180:0:1 and 0.4);
        \draw (-1,0) arc (180:360:1 and 0.4);
        \draw (-1,0) -- (-1,4);
        \draw (1,0) -- (1,4);
        \draw (0,4) ellipse (1 and 0.4);
    \end{scope}
\end{tikzpicture}
\end{minipage}
\par\vspace{0.2cm}
\nmtAnswerBox

\zadtask{Визначте суму всіх цілих значень $a$, за кожного з яких рівняння $$\frac{9 \cos x - 2a - 5}{\sqrt{\pi^2 - 4x^2}} = 0$$ має хоча б один корінь.}
\nmtAnswerBox

\end{document}